\documentclass{article}
\usepackage[utf8]{inputenc}
\usepackage{amsmath,amssymb}
\usepackage[most]{tcolorbox}
\usepackage{geometry}
\geometry{margin=2.5cm}

\newcommand{\step}[1]{\par\noindent\hangindent=3.5em\hangafter=1%
  \makebox[3.5em][l]{\textbf{#1.}}}
\newcommand{\steptag}[1]{\hfill{\small\textsf{[#1]}}}

\newtcolorbox{proofbox}[1]{
  colback=gray!5, colframe=gray!60,
  fonttitle=\bfseries, title={#1},
  breakable, enhanced,
  left=4pt, right=4pt, top=4pt, bottom=4pt,
  before skip=10pt, after skip=10pt
}

\begin{document}

\begin{proofbox}{Route F F2 positive-unital compression: let K >= 1 be a d...}
\step{1} Route F F2 positive-unital compression: let K >= 1 be a dimension-independent constant, n >= 1, Q: l\_inf\textasciicircum{}n -> l\_inf\textasciicircum{}n row-stochastic, D: M\_n -> C\textasciicircum{}n diagonal extraction onto the complex diagonal algebra C\textasciicircum{}n = l\_inf\textasciicircum{}n(C), J: C\textasciicircum{}n -> M\_n diagonal inclusion, Q\_C: C\textasciicircum{}n -> C\textasciicircum{}n the canonical complex-linear extension of Q, and Phi = J Q\_C D, B a finite-dimensional unital C*-algebra, and Delta: B -> M\_n, Upsilon: M\_n -> B UCP maps; if 0 <= eta <= min\{(24K)\textasciicircum{}\{-1\},1\}, ||Delta Upsilon - Phi||\_cb <= K*eta, ||Upsilon Delta - I\_B||\_cb <= K*eta, and ||Upsilon(Delta x Delta y) - xy|| <= K*eta*||x||*||y|| for all x,y in B, then B is commutative and there are k >= 1 and a unital *-isomorphism iota\_C: C\textasciicircum{}k = l\_inf\textasciicircum{}k(C) -> B such that D Delta iota\_C maps R\textasciicircum{}k into R\textasciicircum{}n, iota\_C\textasciicircum{}\{-1\} Upsilon J maps R\textasciicircum{}n into R\textasciicircum{}k, and the resulting restrictions and corestrictions A := (D Delta iota\_C)|\_\{R\textasciicircum{}k\}: l\_inf\textasciicircum{}k -> l\_inf\textasciicircum{}n and M := (iota\_C\textasciicircum{}\{-1\} Upsilon J)|\_\{R\textasciicircum{}n\}: l\_inf\textasciicircum{}n -> l\_inf\textasciicircum{}k are positive unital maps satisfying ||Q - AM|$|\_\{inf-\rangle$inf\} <= K*eta, ||QA - A|$|\_\{inf-\rangle$inf\} <= 2K*eta, and ||Ax||\_inf >= (1-3K*eta)*||x||\_inf for every x in l\_inf\textasciicircum{}k. \steptag{validated} \\[6pt]
\step{1.1} UCP complete contractivity (shared lemma): for unital C*-algebras A,C and every UCP map T:A->C in the sense of def-ucp-map, each amplification T\_r=id\_\{M\_r\} tensor T is contractive, hence ||T||\_cb<=1. In particular ||Delta||\_cb<=1 and ||Upsilon||\_cb<=1. \steptag{validated} \\[6pt]
\step{1.1.1} Matrix Schwarz step from 2-positivity: fix r>=1 and set S=id\_\{M\_r\} tensor T. Complete positivity and unitality make S unital and 2-positive. For a in M\_r(A), the 2-by-2 block [a,1]\textasciicircum{}*[a,1]=[[a* a,a*],[a,1]] is positive in M\_2(M\_r(A)); applying id\_\{M\_2\} tensor S gives [[S(a* a),S(a)*],[S(a),1]]>=0. The Schur complement of the lower-right identity, equivalently positivity tested on columns of the form (xi,-S(a)xi), yields S(a)*S(a)<=S(a* a). \steptag{validated} \\[4pt]
\step{1.1.2} Norm conclusion: positivity and unitality give 0<=S(a* a)<=||a||\textasciicircum{}2 1. Together with node 1.1.1 this yields ||S(a)||\textasciicircum{}2=||S(a)*S(a)||<=||a||\textasciicircum{}2. Hence ||S||<=1 for every r; equality holds on the identity, so ||S||=1 and ||T||\_cb=1. Applying this to the root UCP maps gives ||Delta||\_cb,||Upsilon||\_cb<=1. \steptag{validated} \\[4pt]
\step{1.2} Typed diagonal-range facts: by def-stochastic, Q\_C:C\textasciicircum{}n->C\textasciicircum{}n is positive and unital; since its domain is commutative it is UCP. The diagonal inclusion J is a unital *-homomorphism and hence UCP, while D([a\_ij])=(a\_11,...,a\_nn) is UCP because every amplification sends a positive block matrix [A\_ij] to the direct sum of its positive diagonal compressions. Thus Phi=J Q\_C D is UCP and, by node 1.1, contractive; Phi(a) lies in J(C\textasciicircum{}n), any two elements in its range commute, D J=I\_\{C\textasciicircum{}n\}, and J is isometric. Moreover every positive linear map is *-preserving, so D,Delta,Upsilon,J and iota\_C whenever defined carry self-adjoint elements to self-adjoint elements. \steptag{validated} \\[4pt]
\step{1.3} Approximate invariance: under the root hypotheses, for every b in B one has ||Delta b-Phi(Delta b)|| <= 2K*eta*||b||. Indeed Delta b-Phi Delta b = Delta(b-Upsilon Delta b)+(Delta Upsilon-Phi)(Delta b), whose two terms have norms at most K*eta*||b|| and K*eta*||Delta b||<=K*eta*||b|| by node 1.1. \steptag{validated} \\[4pt]
\step{1.4} Approximate commutativity: under the root hypotheses, for all x,y in B, ||xy-yx|| <= 10K*eta*||x||*||y||. Put a=Delta x, b=Delta y, a0=Phi a, b0=Phi b. By node 1.3, ||a-a0||<=2K*eta||x|| and ||b-b0||<=2K*eta||y||; by nodes 1.1-1.2, ||a||,||a0||<=||x|| and ||b||,||b0||<=||y||, while [a0,b0]=0. Hence ||[a,b]||<=||[a-a0,b]||+||[a0,b-b0]||<=8K*eta||x||||y||. The two approximate-product hypotheses for (x,y) and (y,x), together with contractivity of Upsilon, give ||[x,y]||<=2K*eta||x||||y||+||Upsilon([a,b])||<=10K*eta||x||||y||. \steptag{validated} \\[4pt]
\step{1.5} Norm-two witness lemma, proved without finite-dimensional C*-classification: if a finite-dimensional unital C*-algebra B is noncommutative, then there exist contractions x,y in B with ||xy-yx||=2. \steptag{validated} \\[6pt]
\step{1.5.1} Minimal-projection decomposition: every nonzero finite-dimensional unital C*-algebra B has a finite family of mutually orthogonal minimal projections p\_1,...,p\_m summing to 1. Indeed choose a nonzero projection minimal under subprojection, repeat in the complementary corner, and termination follows because nonzero orthogonal projections are linearly independent. For a minimal projection p, pBp=Cp: otherwise a non-scalar self-adjoint element of pBp has, by finite-dimensional spectral functional calculus, a nonzero proper spectral projection below p. \steptag{validated} \\[4pt]
\step{1.5.2} Off-diagonal corner: for the decomposition in node 1.5.1, noncommutativity implies p\_i B p\_j is nonzero for some i!=j. For if every off-diagonal corner vanished, then b=sum\_\{i,j\}p\_i b p\_j=sum\_i lambda\_i(b)p\_i for every b, using p\_i B p\_i=Cp\_i, so B would be contained in the commutative span of the p\_i. \steptag{validated} \\[4pt]
\step{1.5.3} Matrix-unit normalization: choose distinct minimal p,q and 0!=z in pBq as in node 1.5.2. Since z* z belongs to qBq=Cq and zz* belongs to pBp=Cp, the C*-identity gives z*z=||z||\textasciicircum{}2 q and zz*=||z||\textasciicircum{}2 p. Thus v=z/||z|| obeys v*v=q and vv*=p; also pq=0 and v=pvq imply v\textasciicircum{}2=(v*)\textasciicircum{}2=0. \steptag{validated} \\[4pt]
\step{1.5.4} Pauli witness: set x=p-q and y=v+v*. Then x*=x, y*=y, x*x=y*y=p+q, so ||x||=||y||=1. Direct multiplication using v=pvq, v*v=q and vv*=p gives xy-yx=2(v-v*), while (v-v*)*(v-v*)=p+q has norm 1. Hence ||xy-yx||=2, proving node 1.5. \steptag{validated} \\[4pt]
\step{1.6} Commutativity forcing with constants bound locally: let K>=1 and 0<=eta<=(24K)\textasciicircum{}\{-1\}. If B satisfies node 1.4, then B is commutative. Otherwise node 1.5 supplies contractions x,y with commutator norm 2, whereas node 1.4 gives 2<=10K*eta<=10/24=5/12<2, a contradiction. \steptag{validated} \\[4pt]
\step{1.7} Commutative coordinates: from node 1.6 and external projection-basis-kitaev-1361, there are k>=1 and pairwise orthogonal nonzero projections Pi\_1,...,Pi\_k in B with sum 1 and spanning B. The formula iota\_C:C\textasciicircum{}k->B, iota\_C(z\_1,...,z\_k)=sum\_j z\_j Pi\_j, is a unital *-isomorphism. \steptag{validated} \\[6pt]
\step{1.7.1} Projection-basis input: node 1.6 makes B a finite-dimensional commutative unital C*-algebra. The permitted byte-verbatim external projection-basis-kitaev-1361 states that such an algebra is described by a projection basis \{Pi\_1,...,Pi\_k\} with Pi\_j*=Pi\_j, Pi\_j Pi\_l=delta\_\{jl\}Pi\_l, and sum\_j Pi\_j=1. Because it is a basis, k>=1, every Pi\_j is nonzero, and the Pi\_j linearly span B. \steptag{validated} \\[4pt]
\step{1.7.2} Coordinate calculation: define iota\_C(z)=sum\_j z\_j Pi\_j. The relations in node 1.7.1 give iota\_C(1,...,1)=1, iota\_C(zw)=iota\_C(z)iota\_C(w), and iota\_C(conjugate(z))=iota\_C(z)*. Linear independence makes iota\_C injective and spanning makes it surjective. Hence it is a unital *-isomorphism C\textasciicircum{}k->B. \steptag{validated} \\[4pt]
\step{1.8} Typed positive-unital real maps: for iota\_C from node 1.7, the compositions D Delta iota\_C:C\textasciicircum{}k->C\textasciicircum{}n and iota\_C\textasciicircum{}\{-1\} Upsilon J:C\textasciicircum{}n->C\textasciicircum{}k are UCP, hence positive and unital. Since positive maps are *-preserving, they map the self-adjoint parts R\textasciicircum{}k into R\textasciicircum{}n and R\textasciicircum{}n into R\textasciicircum{}k. Their restrictions/corestrictions A:l\_inf\textasciicircum{}k->l\_inf\textasciicircum{}n and M:l\_inf\textasciicircum{}n->l\_inf\textasciicircum{}k are therefore positive unital real maps. \steptag{validated} \\[4pt]
\step{1.9} AM estimate: for z in R\textasciicircum{}n, use D J=I and Phi=J Q\_C D to identify Qz=D Phi Jz, while AMz=D Delta Upsilon Jz. Since D and J are contractions, ||(Q-AM)z||\_inf <= ||(Phi-Delta Upsilon)Jz|| <= K*eta||z||\_inf. Thus ||Q-AM|$|\_\{inf-\rangle$inf\}<=K*eta. \steptag{validated} \\[4pt]
\step{1.10} QA-A estimate: for x in R\textasciicircum{}k, QAx-Ax=D(Phi Delta iota\_C x-Delta iota\_C x), because D Phi=Q\_C D and Q\_C restricts to Q on R\textasciicircum{}n. Node 1.3 and contractivity of D and the isometry iota\_C yield ||QAx-Ax||\_inf<=2K*eta||x||\_inf. Hence ||QA-A|$|\_\{inf-\rangle$inf\}<=2K*eta. \steptag{validated} \\[4pt]
\step{1.11} Lower-modulus inference only: for x in R\textasciicircum{}k, node 1.1 and ||Upsilon Delta-I\_B||\_cb<=K*eta imply ||Delta iota\_C x|| >= ||Upsilon Delta iota\_C x|| >= (1-K*eta)||x||\_inf. Node 1.3 gives ||Delta iota\_C x-Phi Delta iota\_C x||<=2K*eta||x||\_inf, and ||Phi Delta iota\_C x||=||J Q\_C A x||<=||A x||\_inf by contractivity of Q\_C and isometry of J. Therefore ||A x||\_inf >= (1-3K*eta)||x||\_inf. \steptag{validated} \\[4pt]
\step{1.12} Assembly: nodes 1.6-1.11 provide B commutative, k>=1, the required unital *-isomorphism iota\_C, the stated real-part mapping properties and positive unital maps A,M, together with all three inequalities ||Q-AM|$|\_\{inf-\rangle$inf\}<=K*eta, ||QA-A|$|\_\{inf-\rangle$inf\}<=2K*eta, and ||Ax|$|\_inf\rangle$=(1-3K*eta)||x||\_inf for every x. This is exactly the conclusion of root node 1. \steptag{validated} \\[6pt]
\step{1.12.1} Endpoint bridge for the assembly: node 1.6 is stated for K>=1 and 0<=eta<=(24K)\textasciicircum{}\{-1\}, so it covers the root endpoint eta=0. Concretely, at eta=0 node 1.4 gives ||xy-yx||<=10K*0*||x||*||y||=0 for every x,y in B, hence B is commutative directly; equivalently, if B were noncommutative then node 1.5 and node 1.4 would give 2=||xy-yx||<=0, a contradiction. Thus node 1.6 supplies commutativity on the entire root interval 0<=eta<=min\{(24K)\textasciicircum{}\{-1\},1\}; node 1.7 consequently supplies iota\_C also at eta=0, and nodes 1.8-1.11 then supply the mapping properties and estimates there (with right sides K*eta, 2K*eta, and 1-3K*eta evaluated at eta=0). Therefore node 1.12 assembles the root conclusion without an endpoint gap. \steptag{validated} \\[4pt]
\end{proofbox}

\end{document}
